\subsection{CASIMIR ELEMENT}

PROPOSITION 11. Let \(\mathfrak{g}\) be a Lie algebra over a field K, U its enveloping algebra, \(\mathfrak{h}\) a finite-dimensional ideal of \(\mathfrak{g}\) and \(\beta\) an invariant bilinear form on \(\mathfrak{g}\) , whose restriction to \(\mathfrak{h}\) is non-degenerate. Let \((e_i)_{1\leqslant i\leqslant n}, (e_j')_{1\leqslant j\leqslant n}\) be two bases of \(\mathfrak{h}\) such that \(\beta(e_i,e_j') = \delta_{ij}\) .

Then the element \(c = \sum_{i=1} e_i e_i'\) of U belongs to the centre of U and is independent of the choice of basis \((e_i)\) .

For \(x \in \mathfrak{g}\) let \(x_{\mathfrak{h}}\) be the restriction to \(\mathfrak{h}\) of \(\operatorname{ad}_{\mathfrak{g}} x\) . Then \(x \mapsto x_{\mathfrak{h}}\) is a representation of \(\mathfrak{g}\) on the vector space \(\mathfrak{h}\) and the restriction \(\beta'\) of \(\beta\) to \(\mathfrak{h}\) is invariant under this

\[
\sum_ {i = 1} ^ {n} e _ {i} \otimes e _ {i} ^ {\prime}
\]

representation. By no. 5, Example 3, the tensor \(\sum_{i=1}^{n} e_i \otimes e_i'\) is independent of the choice of basis \((e_i)\) and is an invariant element of the tensor algebra of \(\mathfrak{h}\) . It is also an element of the tensor algebra \(T\) of \(\mathfrak{g}\) , which is invariant for the representation derived from the adjoint representation of \(\mathfrak{g}\) . Its canonical image in U, that is \(c\) , is therefore independent of the choice of basis \((e_i)\) and is an invariant for the representation of \(\mathfrak{g}\) on U considered at the end of no. 2. This element is therefore permutable with every element of \(\mathfrak{g}\) and therefore belongs to the centre of U.

When \(\beta\) is the bilinear form associated with a g-module M, the element \(c\) of Proposition 11 is called the Casimir element associated with M (or with the corresponding representation). This element exists if the restriction of \(\beta\) to \(\mathfrak{h}\) is nondegenerate.

PROPOSITION 12. Let g be a Lie algebra over a field K, h an ideal of g of finite dimension \(n\) and \(\mathbf{M}\) a g-module of finite dimension over \(\mathbf{K}\) . Let \(c\) be the Casimir element (assumed to exist) associated with \(\mathbf{M}\) and \(\mathfrak{h}\) .

\begin{enumerate}
\def\labelenumi{(\alph{enumi})}
\item
  \(\mathrm{Tr}(c_{\mathrm{M}})=n.\)
\item
  If \(\mathbf{M}\) is simple and \(n\) is not divisible by the characteristic of \(\mathbf{K}\) , \(c_{\mathbf{M}}\) is an automorphism of \(\mathbf{M}\) .
\end{enumerate}

In the notation of Proposition 11,

\[
\mathrm{Tr} (c _ {\mathrm{M}}) = \sum_ {i = 1} ^ {n} \mathrm{Tr} ((e _ {i}) _ {\mathrm{M}} (e _ {i} ^ {\prime}) _ {\mathrm{M}}) = \sum_ {i = 1} ^ {n} \beta (e _ {i}, e _ {i} ^ {\prime}) = n.
\]

Hence, if \(n\) is not divisible by the characteristic of \(\mathbf{K}\) , \(c_{\mathbf{M}} \neq 0\) . On the other hand, as \(c\) belongs to the centre of \(\mathrm{U}\) , \(c_{\mathbf{M}}\) is permutable with all the \(x_{\mathbf{M}}\) , \(x \in \mathfrak{g}\) . If further \(\mathbf{M}\) is simple, \(c_{\mathbf{M}}\) is therefore invertible in \(\mathcal{L}(\mathbf{M})\) (Algebra, Chapter VIII, § 4, no. 3, Proposition 2).

